\documentclass[10pt]{article}

\usepackage[letterpaper,margin=1in]{geometry}

\usepackage[T1]{fontenc}
\usepackage{xcolor}
\usepackage{hyperref}

\definecolor{AIDaRBlue}{HTML}{155E75}
\hypersetup{colorlinks=true,urlcolor=AIDaRBlue,linkcolor=AIDaRBlue}
\setcounter{secnumdepth}{0}

\newcommand{\systemurl}{\href{https://submityour.work}{submityour.work}}
\newcommand{\workshopurl}{\href{https://aidar-workshop.github.io/2026/}{AIDaR workshop}}

\title{\vspace{-3em}NeurIPS 2026 AIDaR Workshop: Submission and Review}
\author{A. Sina Booeshaghi}
\date{30 July 2026}

\begin{document}

\maketitle
\begin{center}
\fcolorbox{AIDaRBlue}{AIDaRBlue!6}{%
\parbox{0.94\textwidth}{\small
\textbf{Summary.} The NeurIPS 2026 AIDaR workshop asks how scientific data
should be organized so models and agents can use it. To demonstrate this idea,
we built the AI Data Readiness System (AIDaRS), a simple system for submitting
and reviewing workshop projects. AIDaRS makes the workshop itself a practical
test of AI-data-ready scientific review.}}
\end{center}

\section{Purpose}

The \workshopurl{} is a NeurIPS 2026 workshop on data infrastructure and
benchmarks for reliable scientific AI systems. The workshop brings together
scientists and industry professionals who build scientific data systems and
evaluate AI systems on scientific tasks. It asks how a scientific data
ecosystem should be organized for models and agents, and how this organization
affects how we do science.

The same question applies to workshop submissions. A workshop about AI data
readiness should make its own submissions ready for people, models, and agents.
A PDF submission is important, but it is not the complete scientific record.
Scientific review also requires code, data notes, analysis outputs, and a clear
record of changes that a PDF often omits.

We built AIDaRS, a machine-readable submission system that operates in parallel
with OpenReview. AIDaRS connects three parts of the review process: the author's
work; a submission system that includes a command-line tool, an AI skill, and a
GitHub service hosted on \systemurl{}; and private GitHub repositories where
reviewers inspect projects and record their reviews.

\section{Submission}

The AIDaRS submission process is anonymous and straightforward. Once an author
submits their work through OpenReview, they or an AI agent runs the AIDaRS
command-line tool or AI skill. The operator points to a project directory and
supplies an OpenReview link. The tool prepares a private copy,
removes Git history, performs common privacy-preserving checks, and
reports files that require the author's attention before submission.

The tool then submits the project through the AIDaRS service on \systemurl{}. The service
runs the AIDaRS GitHub App, which creates an anonymous private repository in the
\texttt{NeurIPS2026-AIDaR} organization and opens a pull request for review.
The AIDaRS bot adds the project and later author responses to GitHub. The
author does not need a GitHub account or token. This process keeps the author's
GitHub identity private.

The same AIDaRS tool supports versioned submission updates. An author can change the local project and submit a new version to the existing pull request. Review comments and project history stay in one place while the repository shows the current submission.
Importantly, AIDaRS provides automatic checks for author anonymity, but---like
OpenReview---it cannot guarantee the removal of all identity information from a
project. Authors must inspect their final files before submission.

\section{Review}

The AIDaRS system lets reviewers inspect submissions and make specific scientific
critiques through GitHub. Every reviewer uses a named GitHub account and joins
the \texttt{NeurIPS2026-AIDaR} organization. Reviewers can inspect all private
workshop submissions, and the organizing committee assigns reviewers to each
project. Reviewers can also clone a project when they need to run or examine it
locally.

Reviewers work through familiar GitHub tools. They can submit a formal review,
write a general comment, or comment on a specific line. Submitting authors retrieve these
reviews through the same CLI or skill that they used to submit. An author can
respond without entering GitHub; the AIDaRS bot posts the response to the pull
request and preserves author anonymity.

The GitHub workspace also makes controlled AI evaluation possible. The
organizing committee can run trusted AI evaluation tools on a project and ask
them to review its contents. This creates a practical path for testing whether a
submission is organized well enough for models and agents to use.

\section{Worked example}

The live demonstration uses the private repository
\href{https://github.com/NeurIPS2026-AIDaR/submission-fc04abb65d1e}{%
\texttt{NeurIPS2026-AIDaR/submission-fc04abb65d1e}} and
\href{https://github.com/NeurIPS2026-AIDaR/submission-fc04abb65d1e/pull/1}{its first pull request}.
The example follows five steps:
\begin{enumerate}
  \item The author submits a paper to OpenReview and receives an OpenReview
  link.
  \item A person or their AI agent runs the AIDaRS CLI or skill with the project
  directory and OpenReview link.
  \item AIDaRS checks the project and creates the anonymous private repository
  and pull request in the workshop organization.
  \item The organizing committee assigns reviewers. Reviewers inspect the
  repository and write named comments on the pull request.
  \item The author retrieves the comments, updates the local project, and uses
  AIDaRS to post a response and revision to the same pull request.
\end{enumerate}

\section{Conclusion}

The AIDaRS review system connects infrastructure, tools, and review in one
practical workflow. A person or AI agent uses the same simple tool to submit a
project through \systemurl{}. The service turns the project into an anonymous,
versioned GitHub submission. Reviewers use their normal GitHub identities and
tools, while authors respond through AIDaRS. The system operates alongside
OpenReview.

This infrastructure makes the workshop itself a proof of its central claim.
The workshop asks authors to explain how scientific work should be organized
for models and agents. AIDaRS asks them to submit that work in a form that people,
models, and agents can inspect during review. The resulting AIDaRS system shows
the practical value of AI-data-ready research projects and can support a
post-workshop study of future conference organization.

\end{document}
